\section{Identification Strategy} \label{sec:identification}

To estimate the causal impacts of the program on the outcomes of interest, we estimate the following reduced-form regression:
\begin{align} 
     Y_{iksm}= \alpha + \delta T_{iksm} + \sum_{n=1}^{6} \gamma_n  + \epsilon_{iksm} 
     \label{eq:Y}
\end {align}

\noindent

in which $Y$ is the outcome variable (e.g., financial proficiency) of student $i$ in grade $k$ at school $s$ in municipality $m$, $T$ is a binary variable that is equal to 1 if the student belongs to a school assigned to the program (treatment) and 0 otherwise, $\gamma_n$ are strata fixed effects, and $\epsilon$ is the idiosyncratic error.\footnote{\autoref{tab:sample} shows that the cluster-randomized trial has six clusters.}

The parameter of interest, $\delta$, measures the intention-to-treat (ITT), the treatment effect on students of schools randomly assigned to receive the treatment. The standard error is clustered at the school level, the randomization unit. It is worth noticing that this regression does not distinguish the effect between grades. To assess the program's effect only for elementary, only for middle school, and for each grade included in the pilot, we separately re-estimate \autoref{eq:Y} for each one of these samples.\footnote {We also estimate standard errors clustering at the class level. The results are very similar and are reported in the \onlineappendix.} In addition to that, we follow \cite{firpo2009unconditional} and estimate unconditional quantile intention-to-treat (UQITT) effects to check whether the program has distinct impacts across the financial proficiency distribution.\footnote{We estimate the UQITT in two steps. We first regress the outcome on a constant and the strata dummies and save the residuals; we then use the residuals as the dependent variable in a quantile regression on the treatment dummy. Since the estimates do not control for covariates, we call these unconditional quantile intention-to-treat effects.}

Because each table reports several outcomes at once, we adjust inference for multiple hypothesis testing with Romano--Wolf stepdown p-values \citep{romano2005stepwise, clarke2020romano}, which control the probability of falsely rejecting at least one true null hypothesis when the outcomes of a table are tested jointly. While applications in this literature typically define families of related outcomes within a subsample \citep{frisancho2023school}, we take a more conservative approach and treat all outcome-by-subsample hypotheses reported in each table as a single family, computing the adjusted p-values with 1,000 bootstrap replications that resample school clusters within randomization strata. We report the significance level of every estimate according to this p-value, so the stars in the tables already reflect the multiple-hypothesis correction. The tables also report randomization-inference p-values \citep{young2019channeling}. 


